\chapter{Circuit lower bounds and proof complexity}
\label{chap:lowerbounds}

This chapter collects unconditional lower bounds for concrete computational
models, together with the proof-complexity side of the same program. The
library proves Shannon's counting bounds, which determine the maximum circuit
complexity of an $N$-bit function up to a constant factor, the essential-input
bound and Schnorr's $2N - 1$ bound for parity, Valiant's depth-reduction lemma,
Spira's formula balancing, the monotone Karchmer--Wigderson correspondence, and
the soundness of resolution. For $\cc{AC}^0$ it proves a finite, width-sensitive
switching lemma, iterates it through arbitrary depth-bounded formulas, and
derives a finite counting obstruction for every formula computing parity. All
of these results are finite or nonuniform. The next target is the family-level
theorem that parity is not in $\cc{AC}^0$, which needs only arithmetic on top
of the finite obstruction (roadmap track L4). Beyond it, the plan covers
monotone circuit lower bounds, pseudorandom restrictions, and resolution size
and width lower bounds leading to the pigeonhole principle and to Cook--Reckhow
proof systems (roadmap track L5).
Section~\ref{sec:lowerbounds-algebraic} collects the results of the imported
algebraic-circuits library, stated for CSLib's circuit model: parity bounds in
the De Morgan basis and in $\cc{AC}^0$, monotone CLIQUE, Nechiporuk's formula
bound, Karchmer--Wigderson, a circuit-size hierarchy, conditional complexity,
algebraic lower bounds, and mass production.

\section{Counting bounds}
\label{sec:lowerbounds-counting}

\begin{theorem}[Shannon lower bound]
  \label{thm:lowerbounds-shannon-lower}
  \lean{Complexity.shannon_lower_bound_circuit, Complexity.shannon_sizeComplexity}
  \leanok
  \uses{def:circuit, def:circuits-bases, def:circuits-size-complexity}
  For every $N \ge 6$ there is a function $f : \{0,1\}^N \to \{0,1\}$ that no
  fan-in-two AND/OR circuit of size at most $\lfloor 2^N / (5N) \rfloor$
  computes. Equivalently, $\mathrm{size}(f) > \lfloor 2^N / (5N) \rfloor$.
\end{theorem}
\begin{proof}
  \leanok
  Encode circuits of a given size by short descriptors and compare their
  number with the $2^{2^N}$ functions.
\end{proof}

\begin{theorem}[Shannon upper bound]
  \label{thm:lowerbounds-shannon-upper}
  \lean{Complexity.shannon_upper_bound}
  \leanok
  \uses{def:circuits-size-complexity}
  For every $N \ge 16$, every $f : \{0,1\}^N \to \{0,1\}$ satisfies
  $\mathrm{size}(f) \le \lfloor 18 \cdot 2^N / N \rfloor$ over the fan-in-two
  AND/OR basis.
\end{theorem}
\begin{proof}
  \leanok
  Select among the columns of the truth table with a multiplexer, reading each
  column from a library that contains every function of the remaining
  variables.
\end{proof}

\begin{theorem}[Lupanov]
  \label{thm:lowerbounds-lupanov}
  \lean{Complexity.lupanov_sizeComplexity}
  \leanok
  \uses{def:circuits-size-complexity, thm:circuits-cslib}
  For every $\varepsilon > 0$ and all sufficiently large $N$, every
  $f : \{0,1\}^N \to \{0,1\}$ satisfies
  $\mathrm{size}(f) \le (1 + \varepsilon)\, 2^N / N$.
\end{theorem}
\begin{proof}
  \leanok
  \uses{thm:circuits-cslib}
  CSLib proves the bound for its De Morgan circuits by refining the column
  library: the rows of the truth table are grouped into strips, and each
  strip's column patterns are computed once and shared. A CSLib circuit with
  $g$ gates is a fan-in-two AND/OR circuit of size $g + 1$, and the extra gate
  is absorbed by the slack in $\varepsilon$.
\end{proof}

\begin{theorem}[Shannon lower bound via CSLib]
  \label{thm:lowerbounds-shannon-cslib}
  \lean{Complexity.exists_sizeComplexity_gt_cslib}
  \leanok
  \uses{def:circuits-size-complexity, thm:circuits-cslib}
  For all sufficiently large $N$, some $f : \{0,1\}^N \to \{0,1\}$ satisfies
  $\mathrm{size}(f) > (2^N / N - N) / 2$.
\end{theorem}
\begin{proof}
  \leanok
  \uses{thm:circuits-cslib}
  CSLib's counting argument gives a function that no De Morgan circuit with at
  most $2^N / N$ gates computes, and a fan-in-two AND/OR circuit of size $s$ is
  a CSLib circuit with at most $N + 2s$ gates.
\end{proof}

\begin{theorem}[Circuit size of languages]
  \label{thm:lowerbounds-size-classes}
  \lean{Complexity.mem_SIZE_iff_sliceSizeComplexity_le,
    Complexity.exists_cslib_of_mem_SIZE, Complexity.mem_SIZE_of_cslib,
    Complexity.lupanov_sliceSizeComplexity, Complexity.exists_mem_SIZE_bigO_two_pow_div,
    Complexity.exists_language_not_mem_SIZE,
    Complexity.exists_language_not_mem_SIZE_littleO}
  \leanok
  \uses{def:SIZE, thm:lowerbounds-lupanov, thm:lowerbounds-shannon-cslib, thm:circuits-cslib}
  A language is in $\cc{SIZE}(s)$ exactly when the size complexity of each of
  its length slices is at most $s$, and $\cc{SIZE}$ transfers to CSLib's De
  Morgan circuits with the overheads of Theorem~\ref{thm:circuits-cslib}. Every
  language is in $\cc{SIZE}(s)$ for some $s = O(2^n / n)$, indeed with
  $s(n) \le (1 + \varepsilon) 2^n / n$ for large $n$. Conversely, some language
  lies outside $\cc{SIZE}(s)$ whenever $n + 2s(n) \le 2^n / n$ for large $n$,
  in particular whenever $s = o(2^n / n)$.
\end{theorem}
\begin{proof}
  \leanok
  \uses{thm:lowerbounds-lupanov, thm:lowerbounds-shannon-cslib}
  Apply Lupanov's bound to every slice, with finitely many small lengths
  absorbed into the big-O constant. For the lower bound, take a hard function
  at every large length and let the language consist of its accepted strings.
\end{proof}

\section{Results from the algebraic-circuits library}
\label{sec:lowerbounds-algebraic}

These results come from the algebraic-circuits library, imported wholesale
under \texttt{Complexitylib/Algebraic}. They are stated for CSLib's circuit
model and keep the \texttt{Algebraic} namespace; the consolidation plan in
\texttt{ROADMAP.md} relates them to the statements elsewhere in this chapter.

\begin{theorem}[Parity in the De Morgan basis]
  \label{thm:lowerbounds-algebraic-xor}
  \lean{Algebraic.DeMorgan.xor_lowerBound}
  \leanok
  Every De Morgan circuit computing $n$-input parity has at least $3(n - 1)$
  AND and OR gates; constants and negations are free.
\end{theorem}
\begin{proof}
  \leanok
  Gate elimination: fixing a well-chosen input removes at least three AND or
  OR gates while leaving a parity function on one fewer input.
\end{proof}

\begin{theorem}[Parity is not in $\cc{AC}^0$]
  \label{thm:lowerbounds-algebraic-parity-ac0}
  \lean{Algebraic.AC0.parity_not_computable}
  \leanok
  No family of constant-depth, polynomial-size, unbounded fan-in AND/OR
  circuits with negations at the inputs computes parity. This is stated for
  the imported library's own $\cc{AC}^0$ model; deriving
  Theorem~\ref{thm:lowerbounds-parity-not-ac0} for $\cc{AC}^0$ as defined here
  from it is part of the consolidation plan.
\end{theorem}
\begin{proof}
  \leanok
  Håstad's switching lemma: random restrictions collapse each bottom layer
  of such a circuit to a shallow decision tree, reducing the depth, while
  parity restricts to parity on the remaining inputs.
\end{proof}

\begin{theorem}[Monotone CLIQUE]
  \label{thm:lowerbounds-algebraic-clique}
  \lean{Algebraic.Monotone.Clique.Exponential.twoPow_lt_circuitSize}
  \leanok
  On $N = 2^{20t}$ vertices with $t \ge 4$, every monotone circuit computing
  $2^{4t}$-CLIQUE has more than $2^{t \cdot 2^t}$ gates.
\end{theorem}
\begin{proof}
  \leanok
  Razborov's method of approximations, with the sunflower lemma bounding the
  error introduced at each gate.
\end{proof}

\begin{theorem}[Nechiporuk]
  \label{thm:lowerbounds-algebraic-nechiporuk}
  \lean{Algebraic.Nechiporuk.eventually_sq_le_leaves}
  \leanok
  If a family of Boolean functions is $K(n)$-rectangle-free with
  $K(n) \le n^c$ and has at least $2^{n - 2}$ accepting inputs, then for all
  large $n$ every formula over the full binary basis computing its length-$n$
  member has at least $n^2 / (64(c + 3)\log_2 n)$ leaves.
\end{theorem}
\begin{proof}
  \leanok
  Count subfunctions on each block of a partition of the inputs;
  rectangle-freeness forces the maximal number on every block at once.
\end{proof}

\begin{theorem}[Karchmer--Wigderson]
  \label{thm:lowerbounds-algebraic-kw}
  \lean{Algebraic.KW.formulaDepth_eq_protocolDepth, Algebraic.KW.formulaSize_eq_protocolSize}
  \leanok
  The minimum depth and size of a De Morgan formula for $f$ equal the minimum
  depth and number of leaves of a protocol for the Karchmer--Wigderson game of
  $f$.
\end{theorem}
\begin{proof}
  \leanok
  Formulas and protocols translate into each other gate for move: an AND gate
  is a move of the player holding a rejected input, an OR gate a move of the
  player holding an accepted input.
\end{proof}

\begin{theorem}[Polynomial circuit size hierarchy]
  \label{thm:lowerbounds-algebraic-hierarchy}
  \lean{Algebraic.DeMorgan.polynomialSize_ssubset}
  \leanok
  For real exponents $1 \le a < b$, the class of function families with De
  Morgan circuits of size $O(n^a)$ is strictly contained in the class with
  size $O(n^b)$.
\end{theorem}
\begin{proof}
  \leanok
  Change a function of size at most $n^a$ at a few points at a time: each
  point update changes size by $O(n)$, so some intermediate function has size
  between $n^a$ and $n^b$, while Shannon's bound supplies hard functions.
\end{proof}

\begin{lemma}[Hamming-Lipschitz circuit complexity]
  \label{lem:lowerbounds-algebraic-lipschitz}
  \lean{Algebraic.DeMorgan.complexity_dist_le, Algebraic.DeMorgan.exists_complexity_between}
  \leanok
  The De Morgan circuit complexity of $n$-input Boolean functions changes by
  at most $2n$ per changed truth-table entry. Consequently every threshold
  $t \ge 1$ below the complexity of some function is crossed with overshoot at
  most $2n$: some function has complexity in $(t, t + 2n]$.
\end{lemma}
\begin{proof}
  \leanok
  Patching one truth-table entry costs at most $2n$ gates, by combining the old
  circuit with an indicator of the changed point. Walk from a constant function
  to a hard one entry by entry.
\end{proof}

\begin{definition}[Conditional circuit complexity]
  \label{def:lowerbounds-algebraic-conditional}
  \lean{Cslib.Circuits.Circuit.conditionalGateComplexity}
  \leanok
  The conditional complexity of a target given supplied functions is the
  minimum number of gates computing the target from the original inputs and
  the supplied values as free extra inputs, and $\infty$ when no circuit
  computes it.
\end{definition}

\begin{theorem}[Laws of conditional complexity]
  \label{thm:lowerbounds-algebraic-conditional}
  \lean{Cslib.Circuits.Circuit.conditionalGateComplexity_triangle,
    Algebraic.ConditionalComplexity.Linear.conditionalGateComplexity_eq_min_weight}
  \leanok
  \uses{def:lowerbounds-algebraic-conditional}
  Conditional complexity obeys the triangle inequality: computing $h$ given
  $f$ costs at most computing $h$ given $g$ plus computing $g$ given $f$. Over
  $\mathbb{F}_2$, with gates of fan-in at most two, the conditional complexity
  of a nonzero linear form given linear helpers is exactly a minimum weight,
  whatever other operations the basis contains.
\end{theorem}
\begin{proof}
  \leanok
  Compose the two circuits for the triangle inequality. For linear forms, a
  circuit computing the form must combine at least the minimum number of
  input and helper terms, and such a combination is realized directly.
\end{proof}

\begin{theorem}[Hessian rank bounds multiplications]
  \label{thm:lowerbounds-algebraic-hessian}
  \lean{Algebraic.Applications.hessianRank_lowerBound}
  \leanok
  Every arithmetic circuit computing a polynomial $p$ over a field uses at least
  $\lceil \operatorname{rank} H_p(a) / 2 \rceil$ multiplications, where $H_p(a)$ is
  the Hessian of $p$ at any point $a$.
\end{theorem}
\begin{proof}
  \leanok
  Each multiplication gate raises the rank of the second-order part of the
  computed polynomials by at most two, while additions and scalars do not.
\end{proof}

\begin{theorem}[Waring bound for the squarefree monomial]
  \label{thm:lowerbounds-algebraic-waring}
  \lean{Algebraic.Applications.waringSum_lowerBound}
  \leanok
  Over a field of characteristic zero, any expression of $x_1 x_2 \cdots x_{2n}$
  as a sum of scaled $2n$-th powers of linear forms has at least
  $\binom{2n}{n}$ terms.
\end{theorem}
\begin{proof}
  \leanok
  Compare the dimension of the space of partial derivatives of order $n$ of
  each side.
\end{proof}

\begin{theorem}[Sharp mass production]
  \label{thm:lowerbounds-algebraic-mass-production}
  \lean{Algebraic.MassProduction.BlockInduction.exponentialMassProduction,
    Algebraic.MassProduction.Nonuniform.realSharpMassProduction}
  \leanok
  For every rate $0 \le \gamma < 1$ and $\varepsilon > 0$, for all sufficiently
  large $n$, every Boolean function on $n$ inputs can be evaluated on up to
  $2^{\gamma n}$ independent input blocks by a single De Morgan circuit of cost
  at most $(1/(1 - \gamma) + \varepsilon)\, 2^n / n$.
\end{theorem}
\begin{proof}
  \leanok
  A nonuniform construction that routes all copies through shared table
  lookups: sorting networks gather the requests, and each table entry is
  computed once and broadcast to every copy that needs it.
\end{proof}

\begin{theorem}[Cutwidth bound, conditional]
  \label{thm:lowerbounds-algebraic-cutwidth}
  \lean{Algebraic.Cutwidth.eventually_lt_size_of_pathwidthBound}
  \leanok
  Assume that for every $\xi > 0$ all sufficiently large simple 3-regular
  graphs on $h$ vertices have pathwidth at most $(1/6 + \xi)h$. Then every
  rectangle-free Boolean function family with polynomial rectangle threshold
  and at least $2^{n-2}$ accepting inputs needs more than $(4 - \varepsilon)n$
  gates over the full binary basis, for every $\varepsilon > 0$ and all large
  $n$. The pathwidth bound is an explicit hypothesis of the theorem.
\end{theorem}
\begin{proof}
  \leanok
  Order the gates by a path decomposition of the circuit's wiring graph and
  count the rectangles cut at each position.
\end{proof}

\begin{definition}[The KRW conjecture]
  \label{def:lowerbounds-algebraic-krw}
  \lean{Algebraic.KW.KRWDepth, Algebraic.KW.KRWSize}
  \leanok
  \uses{thm:lowerbounds-algebraic-kw}
  The Karchmer--Raz--Wigderson conjecture, stated with explicit slack: the
  formula depth of a block composition $f \diamond g$ is at least the sum of the
  depths of $f$ and $g$ up to an additive slack, and its formula size is at
  least the product of their sizes up to a multiplicative slack. The
  conjecture is open; it is formalized as a statement only.
\end{definition}

\section{Essential inputs and gate elimination}
\label{sec:lowerbounds-gate-elimination}

\begin{definition}[Essential input]
  \label{def:lowerbounds-essential}
  \lean{Complexity.IsEssentialInput, Complexity.essentialInputs}
  \leanok
  \uses{def:bitstring}
  Input $i$ is essential for $f : \{0,1\}^N \to \{0,1\}^M$ if flipping bit $i$
  changes $f(x)$ for some $x$.
\end{definition}

\begin{theorem}[Essential-input bound]
  \label{thm:lowerbounds-essential}
  \lean{Complexity.Circuit.card_essentialInputs_le_totalFanIn,
    Complexity.Circuit.card_essentialInputs_le_mul_size}
  \leanok
  \uses{def:lowerbounds-essential, def:circuit, def:circuits-bases}
  Over any basis, the number of essential inputs of the function computed by a
  circuit is at most the circuit's total fan-in. For circuits over the
  AND/OR basis of fan-in at most $k$, it is at most $k$ times the size.
\end{theorem}
\begin{proof}
  \leanok
  Every essential input is read by some gate input occurrence.
\end{proof}

\begin{definition}[Parity]
  \label{def:lowerbounds-parity}
  \lean{Complexity.Schnorr.xorBool}
  \leanok
  \uses{def:bitstring}
  The parity family $\oplus$ maps $x \in \{0,1\}^N$ to
  $x_0 \oplus \cdots \oplus x_{N-1}$, with value $0$ on the empty string.
\end{definition}

\begin{theorem}[Schnorr]
  \label{thm:lowerbounds-schnorr}
  \lean{Complexity.schnorr_lower_bound_circuit, Complexity.sizeComplexity_xorBool_ge}
  \leanok
  \uses{def:lowerbounds-parity, def:circuits-size-complexity}
  Every fan-in-two AND/OR circuit computing $N$-bit parity or its complement,
  $N \ge 1$, has size at least $2N - 1$. Hence
  $\mathrm{size}(\oplus_N) \ge 2N - 1$.
\end{theorem}
\begin{proof}
  \leanok
  Gate elimination: fixing a well-chosen input removes at least two gates and
  leaves parity or its complement on one fewer input.
\end{proof}

\section{Depth, formulas, and communication}
\label{sec:lowerbounds-depth}

\begin{theorem}[Valiant's depth reduction]
  \label{thm:lowerbounds-valiant}
  \lean{Complexity.Valiant.depth_reduction}
  \leanok
  Let $G$ be a finite acyclic digraph with $S$ edges whose longest path has at
  most $2^k$ vertices, and let $r \le k$. There is a set $F$ of edges with
  $k\,|F| \le r S$ whose removal leaves a digraph whose longest path has at
  most $2^k / 2^r$ vertices.
\end{theorem}
\begin{proof}
  \leanok
  Label vertices by longest-path depth, partition the edges by the first bit
  in which the labels of their endpoints differ, and remove the $r$ smallest
  classes.
\end{proof}

\begin{corollary}[Depth reduction for circuits]
  \label{cor:lowerbounds-valiant-circuit}
  \uses{def:circuit}
  Let $C$ be a circuit whose dependency graph has longest paths of at most
  $2^k$ vertices, and let $r \le k$. Removing a set $F$ of wires with
  $k\,|F| \le r \cdot \mathrm{totalFanIn}(C)$ leaves a dependency graph whose
  longest paths have at most $2^{k-r}$ vertices.
\end{corollary}
\begin{proof}
  \uses{thm:lowerbounds-valiant, lem:circuits-dependency-graph}
  The dependency graph is acyclic and has at most $\mathrm{totalFanIn}(C)$
  edges.
\end{proof}

\begin{theorem}[Spira]
  \label{thm:lowerbounds-spira}
  \lean{Complexity.BoolFormula.exists_spira_balanced}
  \leanok
  \uses{def:circuits-formula}
  Every Boolean formula $\varphi$ of size $s$ has an equivalent formula of
  depth at most $15 + 12 \lceil \log_2 s \rceil$ and size at most
  $80 s^8 - 4$.
\end{theorem}
\begin{proof}
  \leanok
  Find a subformula of between one third and two thirds of the size, and
  recombine the two cases for its value; recurse on both parts.
\end{proof}

\begin{definition}[Monotone formula]
  \label{def:lowerbounds-monotone-formula}
  \lean{Complexity.MonotoneFormula, Complexity.MonotoneFormula.eval,
    Complexity.MonotoneFormula.depth, Complexity.MonotoneFormula.leaves,
    Complexity.IsMonotoneBoolFun}
  \leanok
  \uses{def:bitstring}
  A monotone formula over $N$ variables is a tree built from variables
  $x_0, \dots, x_{N-1}$ and binary conjunction and disjunction, without
  constants or negation. Its leaves are its variable occurrences. A function
  $f$ is monotone if $x \le y$ pointwise and $f(x) = 1$ imply $f(y) = 1$.
\end{definition}

\begin{lemma}[Monotone formulas]
  \label{lem:lowerbounds-monotone-formula}
  \lean{Complexity.MonotoneFormula.eval_monotone,
    Complexity.MonotoneFormula.card_essentialInputs_le_leaves_of_computes}
  \leanok
  \uses{def:lowerbounds-monotone-formula, def:lowerbounds-essential}
  Every monotone formula computes a monotone function, and a monotone formula
  computing $f$ has at least as many leaves as $f$ has essential inputs.
\end{lemma}
\begin{proof}
  \leanok
  Induction on the formula; an essential input must occur as a leaf.
\end{proof}

\begin{definition}[Karchmer--Wigderson protocol]
  \label{def:lowerbounds-kw}
  \lean{Complexity.KarchmerWigderson.Protocol,
    Complexity.KarchmerWigderson.RootProtocol,
    Complexity.KarchmerWigderson.Protocol.depth}
  \leanok
  \uses{def:bitstring}
  A protocol for the monotone Karchmer--Wigderson game of
  $f : \{0,1\}^N \to \{0,1\}$ is a tree indexed by rectangles $X \times Y$ with
  $X \subseteq f^{-1}(1)$ and $Y \subseteq f^{-1}(0)$, starting from
  $f^{-1}(1) \times f^{-1}(0)$. At an internal node, Alice splits $X$ or Bob
  splits $Y$ by an arbitrary predicate. A leaf names a coordinate $i$ with
  $x_i = 1$ for all $x \in X$ and $y_i = 0$ for all $y \in Y$. The depth is the
  height of the tree.
\end{definition}

\begin{theorem}[Monotone Karchmer--Wigderson]
  \label{thm:lowerbounds-kw}
  \lean{Complexity.KarchmerWigderson.Protocol.exists_protocol_of_formula,
    Complexity.KarchmerWigderson.Protocol.exists_formula_of_protocol,
    Complexity.KarchmerWigderson.Protocol.exists_protocol_depth_iff_formula_depth}
  \leanok
  \uses{def:lowerbounds-kw, def:lowerbounds-monotone-formula}
  For every $f : \{0,1\}^N \to \{0,1\}$ and $d \in \mathbb{N}$, the monotone
  Karchmer--Wigderson game of $f$ has a protocol of depth at most $d$ if and
  only if some monotone formula of depth at most $d$ computes $f$.
\end{theorem}
\begin{proof}
  \leanok
  An OR node is an Alice split and an AND node is a Bob split; both
  translations preserve depth exactly.
\end{proof}

\section{Restrictions, normal forms, and decision trees}
\label{sec:lowerbounds-restrictions}

\begin{definition}[Restriction]
  \label{def:lowerbounds-restriction}
  \lean{Complexity.Restriction, Complexity.Restriction.On,
    Complexity.Restriction.applyTo, Complexity.Restriction.comp,
    Complexity.BoolFormula.restrict}
  \leanok
  \uses{def:circuits-formula}
  A restriction assigns each variable either a fixed bit or ``free''; the
  finite-arity version does this for exactly $N$ variables. Restrictions
  compose, the earlier one taking precedence, and act on assignments and on
  formulas by substituting constants for fixed variables.
\end{definition}

\begin{lemma}[Restriction and evaluation]
  \label{lem:lowerbounds-restrict}
  \lean{Complexity.BoolFormula.eval_restrict, Complexity.BoolFormula.restrict_comp,
    Complexity.BoolFormula.restrict_size, Complexity.BoolFormula.restrict_depth}
  \leanok
  \uses{def:lowerbounds-restriction}
  Evaluating $\varphi|_\rho$ at $\alpha$ equals evaluating $\varphi$ at
  $\rho$ applied to $\alpha$. Restricting by $\rho_1$ and then $\rho_2$ equals
  restricting by their composition, and restriction preserves formula size and
  depth.
\end{lemma}
\begin{proof}
  \leanok
  Induction on the formula.
\end{proof}

\begin{definition}[CNF and DNF]
  \label{def:lowerbounds-normal-forms}
  \lean{Complexity.CNF, Complexity.DNF, Complexity.CNF.width, Complexity.DNF.width,
    Complexity.CNF.complexity, Complexity.DNF.complexity}
  \leanok
  \uses{def:bitstring}
  A CNF (resp.\ DNF) over $N$ variables is a list of clauses (resp.\ terms),
  each a list of literals. Its width is the largest number of literals in a
  clause (term), and its complexity is the number of clauses (terms). These
  finite-arity normal forms are distinct from the CNF formulas used for
  satisfiability and resolution.
\end{definition}

\begin{lemma}[Restricting normal forms]
  \label{lem:lowerbounds-normal-form-restrict}
  \lean{Complexity.CNF.restrict, Complexity.CNF.eval_restrict,
    Complexity.CNF.width_restrict_le, Complexity.DNF.restrict,
    Complexity.DNF.eval_restrict, Complexity.DNF.width_restrict_le}
  \leanok
  \uses{def:lowerbounds-normal-forms, def:lowerbounds-restriction}
  Every CNF and DNF has a simplified restriction that computes it under the
  restriction and whose width does not exceed the original width.
\end{lemma}
\begin{proof}
  \leanok
  Delete falsified literals, and delete satisfied clauses (falsified terms).
\end{proof}

\begin{definition}[Decision tree]
  \label{def:lowerbounds-decision-tree}
  \lean{Complexity.DecisionTree, Complexity.DecisionTree.On,
    Complexity.DecisionTree.On.eval, Complexity.DecisionTree.On.depth}
  \leanok
  \uses{def:bitstring}
  A decision tree queries a variable at each internal node and outputs a bit
  at each leaf. The finite-arity version queries only variables below $N$.
  Its depth is the length of a longest root-to-leaf path.
\end{definition}

\begin{lemma}[Decision trees to normal forms]
  \label{lem:lowerbounds-dt-normal-form}
  \lean{Complexity.DecisionTree.On.eval_toDNF, Complexity.DecisionTree.On.eval_toCNF,
    Complexity.DecisionTree.On.width_toDNF_le_depth,
    Complexity.DecisionTree.On.width_toCNF_le_depth}
  \leanok
  \uses{def:lowerbounds-decision-tree, def:lowerbounds-normal-forms}
  A decision tree of depth $s$ is computed by a DNF of width at most $s$ and
  by a CNF of width at most $s$.
\end{lemma}
\begin{proof}
  \leanok
  Take one term per accepting path, and dually one clause per rejecting path.
\end{proof}

\begin{theorem}[Parity at depth two]
  \label{thm:lowerbounds-parity-depth-two}
  \lean{Complexity.DNF.two_pow_le_complexity_of_xorBool,
    Complexity.CNF.two_pow_le_complexity_of_xorBool}
  \leanok
  \uses{def:lowerbounds-normal-forms, def:lowerbounds-parity}
  For $N \ge 1$, every DNF and every CNF computing $N$-bit parity has at
  least $2^{N-1}$ terms, respectively clauses.
\end{theorem}
\begin{proof}
  \leanok
  A satisfiable term that omits a variable accepts two inputs of different
  parity, so every satisfiable term accepts exactly one input, while parity
  accepts $2^{N-1}$ inputs. CNFs follow by negation.
\end{proof}

\section{The switching lemma and $\cc{AC}^0$}
\label{sec:lowerbounds-ac0}

\begin{definition}[Sparse random restriction]
  \label{def:lowerbounds-random-restriction}
  \lean{Complexity.RandomRestriction.Seed, Complexity.RandomRestriction.decode,
    Complexity.RandomRestriction.eventCount}
  \leanok
  \uses{def:lowerbounds-restriction}
  For $q \in \mathbb{N}$, a seed assigns to each of $N$ coordinates one of
  $2q + 1$ symbols: one free symbol and $q$ labelled copies of each fixed bit.
  A seed decodes to a restriction, so a uniform seed leaves each coordinate
  free with probability $1/(2q+1)$ and fixes it to each bit with probability
  $q/(2q+1)$, independently. Events are measured by exact counts of seeds
  among the $(2q + 1)^N$.
\end{definition}

\begin{theorem}[Switching lemma, counting form]
  \label{thm:lowerbounds-switching}
  \lean{Complexity.DNF.switchingBad,
    Complexity.DNF.switchingBad_width_encoding_bound,
    Complexity.DNF.switchingBad_consistentPart_width_encoding_bound,
    Complexity.CNF.switchingBad_width_encoding_bound,
    Complexity.CNF.switchingBad_consistentPart_width_encoding_bound,
    Complexity.DNF.eval_consistentPart, Complexity.CNF.eval_consistentPart}
  \leanok
  \uses{def:lowerbounds-random-restriction, def:lowerbounds-normal-forms,
    def:lowerbounds-decision-tree}
  Let $\varphi$ be a DNF of width $w$ over $N$ variables, no term of which
  contains a complementary pair of literals, and let $q, s \in \mathbb{N}$.
  Call a restriction bad if the canonical decision tree of the restricted
  formula has depth at least $s$. Then
  \[ \#\{\text{bad seeds}\} \cdot q^s \le (2q + 1)^N \, (4(w + 1))^s . \]
  Every DNF agrees with such a sub-DNF of no larger width, and the same bounds
  hold for CNFs. Dividing by $(2q+1)^N q^s$: a uniform seed, which leaves each
  coordinate free with probability $1/(2q+1)$, is bad with probability at most
  $(4(w+1)/q)^s$ when $q \ge 1$.
\end{theorem}
\begin{proof}
  \leanok
  Encode each bad restriction, together with the labels of its first $s$
  queried coordinates, injectively as a seed plus $s$ advice symbols from an
  alphabet of size $4(w + 1)$ (Razborov's encoding). CNFs follow by De Morgan
  duality.
\end{proof}

\begin{definition}[$\cc{AC}^0$ formula]
  \label{def:lowerbounds-ac0-formula}
  \lean{Complexity.AC0Formula, Complexity.AC0Formula.eval, Complexity.AC0Formula.size,
    Complexity.AC0Formula.depth}
  \leanok
  \uses{def:bitstring}
  A negation-normal unbounded formula over $N$ variables is a tree built from
  constants, literals, and AND and OR gates of any fan-in. Its size counts all
  nodes, and its depth counts gates on a longest root-to-leaf path.
\end{definition}

\begin{theorem}[Normalization of $\cc{AC}^0$ circuits]
  \label{thm:lowerbounds-ac0-normalization}
  \lean{Complexity.Circuit.outputAC0Formula, Complexity.Circuit.outputAC0Formula_spec}
  \leanok
  \uses{def:lowerbounds-ac0-formula, def:circuits-bases}
  Let $C$ be an unbounded-fan-in AND/OR circuit with $N$ inputs and $G$
  internal gates, and let output $j$ have depth $D$. Some negation-normal
  unbounded formula computes output $j$, has depth at most $D$, and has size at
  most $(2(N + G) + 1)^{D+1}$.
\end{theorem}
\begin{proof}
  \leanok
  Push negations to the literals by De Morgan duality, remove duplicate signed
  inputs of each gate, and unfold the circuit into a tree.
\end{proof}

\begin{theorem}[Iterated switching]
  \label{thm:lowerbounds-staged-switching}
  \lean{Complexity.AC0Formula.stagedDecisionTree,
    Complexity.AC0Formula.eval_stagedDecisionTree,
    Complexity.AC0Formula.stageEventCount_stagedBad_mul_pow_le,
    Complexity.AC0Formula.exists_shallow_stagedDecisionTree_of_counting}
  \leanok
  \uses{def:lowerbounds-ac0-formula, def:lowerbounds-random-restriction,
    def:lowerbounds-decision-tree}
  Let $\varphi$ be a negation-normal unbounded formula over $N$ variables with
  depth at most $d$, and let $s \ge 2$ and $q \in \mathbb{N}$. Applying $d$
  independent sparse restrictions, one per level, turns $\varphi$ into a
  decision tree that computes $\varphi$ under the composed restriction. The
  number of $d$-tuples of seeds for which this decision tree has depth at least
  $s$, multiplied by $q^s$, is at most
  $\mathrm{size}(\varphi)\,((2q+1)^N)^d\,(4(s + 1))^s$. Consequently, if $q > 0$
  and
  \[ ((2q+1)^N)^d\, s\, q^s + \mathrm{size}(\varphi)\,((2q+1)^N)^d\,(4(s+1))^s N
     < N\, ((2q+1)^{N-1})^d\, q^s , \]
  then some tuple of seeds leaves at least $s$ variables free and reduces
  $\varphi$ to a decision tree of depth less than $s$.
\end{theorem}
\begin{proof}
  \leanok
  \uses{thm:lowerbounds-switching, lem:lowerbounds-dt-normal-form,
    lem:lowerbounds-normal-form-restrict}
  At each level the children are already decision trees of depth below $s$,
  hence CNFs or DNFs of width below $s$, so the switching lemma applies at
  every node; a union bound over the nodes and the exact first moment of the
  number of surviving variables give the criterion.
\end{proof}

\begin{theorem}[Finite parity obstruction]
  \label{thm:lowerbounds-parity-obstruction}
  \lean{Complexity.AC0Formula.parity_counting_obstruction}
  \leanok
  \uses{def:lowerbounds-ac0-formula, def:lowerbounds-parity}
  Let $\varphi$ be a negation-normal unbounded formula over $N$ variables
  that computes parity, and let $d \ge \mathrm{depth}(\varphi)$, $s \ge 2$ and
  $q \ge 1$. Then
  \[ N\, ((2q+1)^{N-1})^d\, q^s \le ((2q+1)^N)^d\, s\, q^s
     + \mathrm{size}(\varphi)\,((2q+1)^N)^d\,(4(s+1))^s N . \]
\end{theorem}
\begin{proof}
  \leanok
  \uses{thm:lowerbounds-staged-switching}
  Parity restricted to a set of free variables needs a decision tree querying
  all of them, so the conclusion of Theorem~\ref{thm:lowerbounds-staged-switching}
  is impossible and its numeric hypothesis must fail.
\end{proof}

\begin{theorem}[Parity is not in $\cc{AC}^0$]
  \label{thm:lowerbounds-parity-not-ac0}
  \uses{def:NC-AC, def:lowerbounds-parity}
  The parity family is not in $\cc{AC}^0$.
\end{theorem}
\begin{proof}
  \uses{thm:lowerbounds-ac0-normalization, thm:lowerbounds-parity-obstruction}
  Normalize a constant-depth polynomial-size family into formulas of constant
  depth $d$ and polynomial size, then choose $q$ and $s$ as powers of $N$ so
  that the inequality of Theorem~\ref{thm:lowerbounds-parity-obstruction}
  fails for large $N$.
\end{proof}

\begin{theorem}[Håstad's size bound]
  \label{thm:lowerbounds-hastad-size}
  \uses{def:lowerbounds-parity, def:circuits-bases}
  For every $d \ge 2$ there is $\varepsilon_d > 0$ such that, for all
  sufficiently large $N$, every unbounded-fan-in AND/OR circuit of depth $d$
  computing $N$-bit parity has size at least $2^{\varepsilon_d N^{1/(d-1)}}$.
\end{theorem}
\begin{proof}
  \uses{thm:lowerbounds-ac0-normalization, thm:lowerbounds-staged-switching}
  The same argument with the restriction parameters optimized, keeping the
  bottom fan-in bound separate from the size.
\end{proof}

\begin{lemma}[Parity in $\cc{TC}^0$]
  \label{lem:lowerbounds-parity-tc0}
  \uses{def:circuits-TC, def:lowerbounds-parity}
  The parity family is in $\cc{TC}^0$.
\end{lemma}
\begin{proof}
  For each odd $k \le N$, the conjunction of the threshold gates
  ``at least $k$ ones'' and ``not at least $k + 1$ ones'' detects exactly $k$
  ones; parity is the disjunction of these, a depth-three circuit of linear
  size.
\end{proof}

\begin{corollary}
  \label{cor:lowerbounds-ac0-ne-tc0}
  \uses{def:NC-AC, def:circuits-TC}
  $\cc{AC}^0 \subsetneq \cc{TC}^0$.
\end{corollary}
\begin{proof}
  \uses{prop:circuits-depth-classes, thm:lowerbounds-parity-not-ac0,
    lem:lowerbounds-parity-tc0}
  $\cc{AC}^0 \subseteq \cc{TC}^0$ is formalized, and parity separates them.
\end{proof}

\begin{theorem}[Pseudorandom switching lemma]
  \label{thm:lowerbounds-pseudorandom-restrictions}
  \uses{thm:lowerbounds-switching}
  \notready
  The bound of Theorem~\ref{thm:lowerbounds-switching} holds, up to an additive
  error $\varepsilon$, when the free/fixed pattern and the fixed values are
  drawn from a distribution that $\varepsilon$-fools CNFs of size
  polynomial in $N$ instead of the uniform product distribution
  (Trevisan--Xue). With an explicit generator fooling CNFs, such restrictions
  can be sampled from $\mathrm{polylog}(N / \varepsilon)$ random bits.
\end{theorem}
\begin{proof}
  The bad event of the switching lemma is itself decided by a small CNF-like
  test on the restriction, so fooling that test preserves its probability.
  This requires a formal notion of pseudorandom distributions and of fooling a
  class of tests, which the library does not yet state; pseudorandom
  generators for $\cc{AC}^0$ would build on it.
\end{proof}

\section{Monotone circuits}
\label{sec:lowerbounds-monotone}

\begin{definition}[Monotone circuit]
  \label{def:lowerbounds-monotone-circuit}
  \uses{def:circuit, def:circuits-bases}
  A monotone circuit is an AND/OR circuit all of whose negation flags are
  false. Every monotone circuit computes a monotone function.
\end{definition}

\begin{theorem}[Monotone Boolean matrix product]
  \label{thm:lowerbounds-monotone-matrix}
  \uses{def:lowerbounds-monotone-circuit}
  \notready
  Every monotone fan-in-two circuit computing the product of two
  $n \times n$ Boolean matrices has at least $n^3$ AND gates.
\end{theorem}
\begin{proof}
  Each product $x_{ik} y_{kj}$ is a prime implicant of output $(i, j)$, and a
  monotone circuit must realize each of the $n^3$ of them at a distinct AND
  gate (Pratt; Paterson; Mehlhorn--Galil). The bound is monotone-specific:
  circuits with negation compute the product with $O(n^{\log_2 7})$ gates
  through integer matrix multiplication.
\end{proof}

\begin{lemma}[Sunflower lemma]
  \label{lem:lowerbounds-sunflower}
  Every family of more than $\ell!\,(p-1)^\ell$ sets, each of size $\ell$,
  contains $p$ sets whose pairwise intersections all equal their common
  intersection.
\end{lemma}
\begin{proof}
  Induction on $\ell$: either a maximal pairwise-disjoint subfamily has $p$
  members, or some element lies in many sets and the lemma applies to them
  with that element removed (Erd\H{o}s--Rado).
\end{proof}

\begin{theorem}[Monotone clique lower bound]
  \label{thm:lowerbounds-monotone-clique}
  \uses{def:lowerbounds-monotone-circuit}
  \notready
  There is $\varepsilon > 0$ such that for all $k \le n^{1/4}$, every monotone
  circuit deciding whether an $n$-vertex graph, given by its adjacency bits,
  contains a $k$-clique has size at least $2^{\varepsilon \sqrt{k}}$
  (Razborov; Alon--Boppana; Arora--Barak Theorem~14.7).
\end{theorem}
\begin{proof}
  \uses{lem:lowerbounds-sunflower}
  Razborov's method of approximations: replace each gate by an approximator
  that is a small union of clique indicators, use the sunflower lemma to keep
  approximators small, and show that each replacement errs on few positive or
  negative test graphs.
\end{proof}

\section{Proof complexity}
\label{sec:lowerbounds-proof-complexity}

\begin{definition}[Resolution]
  \label{def:lowerbounds-resolution}
  \lean{Complexity.SAT.Clause.resolvent, Complexity.SAT.CNF.Derives}
  \leanok
  Clauses are lists of literals over variables in $\mathbb{N}$, and a CNF is a
  list of clauses. The resolvent of clauses $C_1$ and $C_2$ on a variable $v$
  deletes every occurrence of $x_v$ from $C_1$ and of $\neg x_v$ from $C_2$
  and concatenates the rest. A clause is derivable from $\varphi$ if it is a
  clause of $\varphi$ or a resolvent of two derivable clauses. There is no
  weakening rule, and the pivot need not occur in the parents.
\end{definition}

\begin{lemma}[Resolution step]
  \label{lem:lowerbounds-resolvent}
  \lean{Complexity.SAT.Clause.resolvent_sound, Complexity.SAT.Clause.resolvent_length_le}
  \leanok
  \uses{def:lowerbounds-resolution}
  Every assignment satisfying $C_1$ and $C_2$ satisfies their resolvent on any
  variable, and the resolvent has at most $|C_1| + |C_2|$ literals.
\end{lemma}
\begin{proof}
  \leanok
  Case on the value of $x_v$: the parent whose pivot literal is false keeps a
  true literal in the resolvent.
\end{proof}

\begin{theorem}[Soundness of resolution]
  \label{thm:lowerbounds-resolution-sound}
  \lean{Complexity.SAT.CNF.entails_of_derives, Complexity.SAT.CNF.refutation_sound}
  \leanok
  \uses{def:lowerbounds-resolution}
  Every clause derivable from $\varphi$ is satisfied by every assignment
  satisfying $\varphi$. In particular, if the empty clause is derivable from
  $\varphi$, then $\varphi$ is unsatisfiable.
\end{theorem}
\begin{proof}
  \leanok
  \uses{lem:lowerbounds-resolvent}
  Induction on the derivation.
\end{proof}

\begin{theorem}[Completeness of resolution]
  \label{thm:lowerbounds-resolution-complete}
  \uses{def:lowerbounds-resolution}
  If $\varphi$ is unsatisfiable, then the empty clause is derivable from
  $\varphi$.
\end{theorem}
\begin{proof}
  Induction on the number of variables, eliminating one variable by resolving
  all pairs of clauses on it (Davis--Putnam).
\end{proof}

\begin{definition}[Resolution refutations and their measures]
  \label{def:lowerbounds-resolution-proof}
  \uses{def:lowerbounds-resolution}
  A resolution refutation of $\varphi$ is a sequence of clauses ending in the
  empty clause, each a clause of $\varphi$ or a resolvent of two earlier ones.
  Its size is its length and its width is the largest number of literals in
  one of its clauses. A clause is derivable exactly when it ends such a
  sequence.
\end{definition}

\begin{theorem}[Ben-Sasson--Wigderson size--width relation]
  \label{thm:lowerbounds-size-width}
  \uses{def:lowerbounds-resolution-proof}
  \notready
  If an unsatisfiable CNF over $n$ variables whose clauses have at most $k$
  literals has a resolution refutation of size $S$, then it has one of width at
  most $k + O(\sqrt{n \log S})$.
\end{theorem}
\begin{proof}
  Repeatedly restrict a variable occurring in many wide clauses, and combine
  the refutations of the two restrictions.
\end{proof}

\begin{theorem}[Haken]
  \label{thm:lowerbounds-php}
  \uses{def:lowerbounds-resolution-proof}
  \notready
  Let $\mathrm{PHP}^{n+1}_n$ be the CNF over variables $p_{ij}$
  ($i \le n + 1$, $j \le n$) stating that every pigeon $i$ sits in some hole
  and no hole holds two pigeons. Every resolution refutation of
  $\mathrm{PHP}^{n+1}_n$ has size $2^{\Omega(n)}$.
\end{theorem}
\begin{proof}
  \uses{thm:lowerbounds-resolution-sound}
  Restrict the refutation so that every wide clause disappears, then show that
  any refutation of the remaining pigeonhole instance contains a wide clause
  (Beame--Pitassi), or use Haken's bottleneck counting.
\end{proof}

\begin{definition}[Cook--Reckhow proof system]
  \label{def:lowerbounds-proof-system}
  \uses{def:P}
  A propositional proof system for unsatisfiable CNFs is a polynomial-time
  decidable relation $V(\varphi, \pi)$ that is sound ($V(\varphi, \pi)$
  implies that $\varphi$ is unsatisfiable) and complete (every unsatisfiable
  $\varphi$ has some $\pi$ with $V(\varphi, \pi)$). It is polynomially bounded
  if every unsatisfiable $\varphi$ has such a $\pi$ with
  $|\pi| \le p(|\varphi|)$ for a fixed polynomial $p$.
\end{definition}

\begin{proposition}[Resolution is a proof system]
  \label{prop:lowerbounds-resolution-proof-system}
  \uses{def:lowerbounds-proof-system, def:lowerbounds-resolution-proof}
  \notready
  Checking a resolution refutation, encoded as a clause sequence, is a
  Cook--Reckhow proof system.
\end{proposition}
\begin{proof}
  \uses{thm:lowerbounds-resolution-sound, thm:lowerbounds-resolution-complete}
  Checking each step takes polynomial time; soundness and completeness are the
  two theorems above.
\end{proof}

\begin{theorem}[Cook--Reckhow]
  \label{thm:lowerbounds-cook-reckhow}
  \uses{def:lowerbounds-proof-system, def:NP, def:coNP}
  \notready
  A polynomially bounded propositional proof system exists if and only if
  $\cc{NP} = \cc{coNP}$.
\end{theorem}
\begin{proof}
  \uses{thm:cook-levin}
  Unsatisfiability of CNFs is $\cc{coNP}$-complete. A polynomially bounded
  proof system puts it in $\cc{NP}$; conversely an $\cc{NP}$ verifier for it is
  a polynomially bounded proof system.
\end{proof}
